% !TeX spellcheck = ru_RU
% !TEX root = vkr.tex

\section{Эксперимент и апробация}
\label{sec:experiment}

\subsection{Цель и методология}

Целью эксперимента является количественная оценка точности автоматических морфометрических измерений
системы путём сравнения с ручными измерениями, выполненными специалистом лаборатории Гербарий ГБС РАН.

Тестовая выборка включает 9~изображений листьев \textit{Sphagnum} трех разных типов.
Для каждого изображения специалист выполнил ручное измерение двух морфометрических характеристик:
\begin{itemize}
    \item длины листа по главной оси;
    \item поперечных профилей ширины в семи равноотстоящих сечениях
          (на долях 0.125, 0.25, 0.375, 0.5, 0.625, 0.75, 0.875 от полной длины).
\end{itemize}

Те же изображения были обработаны разработанной системой через режим автоматической сегментации и параметрических измерений.
Маски сохранены в директории \texttt{metric\_test/results/};
ручные эталонные измерения~--- в \texttt{metric\_test/mapped/}.

\subsection{Сравнение длин листьев}

В таблице~\ref{tab:lengths} приведены автоматические и ручные измерения длины для всех девяти образцов.
Все автоматические значения превышают ручные, что свидетельствует о систематическом завышении.

\begin{table}[H]
\centering
\caption{Сравнение автоматических и ручных измерений длины листа (пиксели)}
\label{tab:lengths}
\begin{tabular}{crrrr}
\hline
\textbf{№} & \textbf{Авто} & \textbf{Вручную} & $\Delta$, px & APE, \% \\
\hline
1 & 1131.9 & 1106.2 & +25.7  &  2.3 \\
2 &  734.2 &  703.2 & +31.0  &  4.4 \\
3 & 1142.9 & 1132.0 & +10.9  &  1.0 \\
4 & 1511.2 & 1446.2 & +65.0  &  4.5 \\
5 & 2018.4 & 1992.2 & +26.2  &  1.3 \\
6 &  850.9 &  766.9 & +84.0  & 11.0 \\
7 & 1899.2 & 1680.4 & +218.8 & 13.0 \\
8 &  855.2 &  829.2 & +26.0  &  3.1 \\
9 &  443.8 &  408.2 & +35.6  &  8.7 \\
\hline
\textbf{Итого} & & & \textbf{+58.1} & \textbf{5.5} \\
\hline
\end{tabular}
\end{table}

Средняя абсолютная погрешность составила MAE\,=\,58.1\,px при RMSE\,=\,84.0\,px
и средней относительной погрешности MAPE\,=\,5.5\,\%.
Коэффициент корреляции $R = 0.993$ свидетельствует о высокой линейной согласованности
автоматических и ручных измерений.

Наибольшие отклонения наблюдаются у образцов № 7 и № 6
(APE\,=\,13.0\,\% и 11.0\,\%) и объясняются субъективностью ручного измерения:
у этих образцов в основании листа имеются выросты~--- структуры,
относящиеся к листу, которые эксперт при ручной разметке решил не включать в
измеряемую длину. Алгоритм же по полной маске учитывает их как часть листа,
что и даёт расхождение. Для остальных образцов с более регулярной формой
основания относительная погрешность длины не превышает 5\,\%.

\subsection{Сравнение поперечных профилей ширины}

В таблице~\ref{tab:cross} представлены метрики погрешности для поперечных профилей
по каждому из девяти образцов. Анализ проводился по всем 63~значениям (9 листьев $\times$ 7 сечений).

\begin{table}[H]
\centering
\caption{Метрики погрешности поперечных профилей ширины (пиксели)}
\label{tab:cross}
\begin{tabular}{crrrr}
\hline
\textbf{№} & \textbf{Bias} & \textbf{MAE} & \textbf{RMSE} & \textbf{MAPE, \%} \\
\hline
1 & $-3.9$ &  7.7 & 10.7 & 1.95 \\
2 & $-3.9$ &  9.2 & 11.2 & 1.75 \\
3 & $-1.2$ &  2.3 &  2.7 & 0.60 \\
4 &$-12.9$ & 13.7 & 16.3 & 1.54 \\
5 &$-14.3$ & 14.4 & 20.8 & 2.41 \\
6 & $+7.1$ & 13.7 & 16.2 & 3.97 \\
7 &$+12.1$ & 20.7 & 25.3 & 2.44 \\
8 & $-1.8$ &  3.9 &  4.2 & 0.96 \\
9 & $-3.5$ &  4.4 &  7.1 & 1.43 \\
\hline
\textbf{Итого} & $\mathbf{-2.5}$ & \textbf{10.0} & \textbf{14.6} & \textbf{1.90} \\
\hline
\end{tabular}
\end{table}

Суммарная MAPE по поперечным профилям составила \textbf{1.90\,\%}~. Коэффициент корреляции $R = 0.998$ при 63~парах наблюдений подтверждает
высокую согласованность алгоритма сканирования поперечных сечений с ручными измерениями.

Систематический сдвиг по профилям незначителен (Bias\,=\,$-2.5$\,px) и не имеет
выраженного направления: у одних образцов алгоритм слегка занижает ширину (листья 1, 2, 4, 5),
у других~--- слегка завышает (листья 6, 7).

\subsection{Итоговые показатели}

В таблице~\ref{tab:summary} сведены итоговые метрики по двум типам измерений.

\begin{table}[H]
\centering
\small
\caption{Итоговые метрики точности автоматических измерений (Bias, MAE, RMSE~--- в пикселях)}
\label{tab:summary}
\begin{tabular}{lrrrrrr}
\hline
\textbf{Тип измерения} & $N$ & \textbf{Bias} & \textbf{MAE} & \textbf{RMSE} & \textbf{MAPE, \%} & $R$ \\
\hline
Длина листа              &  9 & $+58.1$ & 58.1 & 84.0 & 5.5  & 0.993 \\
Поперечные профили ширины& 63 &  $-2.5$ & 10.0 & 14.6 & 1.90 & 0.998 \\
\hline
\end{tabular}
\end{table}

Измерение поперечных профилей демонстрирует высокую точность~--- погрешность
менее 2\,\% приемлема для сравнительного морфометрического анализа.
Длина листа имеет систематическое завышение порядка 5--6\,\% относительно
ручных измерений: алгоритм по полной маске учитывает выросты у основания
листа, которые эксперт при ручной разметке субъективно исключает (см. анализ
образцов № 7 и № 6 выше).

\subsection{Сравнение по затратам времени}

Среднее время ручного измерения одного листа экспертом составило
$2{,}6 \pm 0{,}4$~минут; автоматическая обработка одного изображения заняла
$\approx 1{,}0$~с (раздел~\ref{sec:auto_perf}), что соответствует ускорению
приблизительно на два порядка.

\subsection{Апробация}

Прототип системы был продемонстрирован консультанту~--- м.н.с.\ лаборатории Гербарий ГБС РАН
Шкурко~А.~В. По итогам демонстрации качество автоматической сегментации признано
удовлетворительным для дальнейшего развития системы.
Видеозапись демонстрации прилагается к работе.

В ходе апробации были выявлены следующие направления доработки:
\begin{itemize}
    \item поддержка масштабного коэффициента (пересчёт пикселей в мм);
    \item возможность экспорта визуализаций в дополнение к CSV;
    \item расстановка точек вдоль периметра с заданным шагом (требование из технического задания).
\end{itemize}
